\section{Pull Request}

La pull request è il meccanismo formale con cui si propone l'integrazione delle modifiche nel branch \texttt{main}. Tutte le modifiche, anche quelle apparentemente piccole, devono passare per una pull request: non sono ammessi push diretti su \texttt{main}.

\subsection{Apertura della Pull Request}

\subsubsection{Step 1: Accedere alla repository su GitHub}

Navigare a \url{https://github.com/Idrorepellente/quantsys_institutional}. Dopo aver eseguito il push del branch, GitHub mostra automaticamente un pulsante \textbf{Compare \& pull request}.

In alternativa, è possibile aprire manualmente la pull request:

\begin{enumerate}
  \item Navigare nella tab \textbf{Pull requests};
  \item Cliccare su \textbf{New pull request};
  \item Selezionare \texttt{main} come \emph{base} e il proprio branch come \emph{compare};
  \item Cliccare su \textbf{Create pull request}.
\end{enumerate}

\subsubsection{Step 2: Compilare i dettagli della Pull Request}

Una buona pull request comunica chiaramente cosa è stato cambiato, perché, e come testarlo.

\paragraph{Titolo} Deve essere descrittivo e conciso. Esempi:
\begin{itemize}
  \item Buono: \emph{Aggiungi exit policy ATR-based al framework strategie}
  \item Cattivo: \emph{Fix}, \emph{Updates}, \emph{Changes}
\end{itemize}

\paragraph{Descrizione} Deve includere:
\begin{itemize}
  \item Riferimento alla issue chiusa: \texttt{Closes \#<numero>};
  \item Descrizione sintetica dei cambiamenti effettuati;
  \item Eventuali note importanti (breaking change, dipendenze nuove, test da eseguire manualmente);
  \item Per modifiche al codice di backtest o alle strategie, includere i numeri OOS pre/post modifica.
\end{itemize}

\subsection{Template di descrizione}

Si raccomanda l'uso del seguente template per le pull request:

\begin{lstlisting}
Closes #<numero-issue>

## Descrizione
<descrizione sintetica dei cambiamenti>

## Cambiamenti principali
- <punto 1>
- <punto 2>
- <punto 3>

## Test eseguiti
- [ ] Test unitari della cartella tests/
- [ ] Walk-forward della strategia X (Sharpe pre: 0.81, post: 0.83)
- [ ] Verifica che la dashboard si avvii correttamente

## Note
<eventuali note: breaking changes, dipendenze nuove,
manual steps richiesti per applicare la modifica>
\end{lstlisting}

\subsection{Code review}

Una volta aperta la pull request:

\begin{enumerate}
  \item Almeno un altro membro del team deve effettuare la review;
  \item Il reviewer può:
  \begin{itemize}
    \item \textbf{Approve}: la PR è pronta per il merge;
    \item \textbf{Request changes}: sono richieste modifiche prima del merge;
    \item \textbf{Comment}: feedback non bloccante.
  \end{itemize}
  \item L'autore della PR deve rispondere ai commenti, applicare le modifiche richieste, e successivamente effettuare nuovo push del branch (la PR si aggiorna automaticamente);
  \item Una volta ottenuta l'approvazione, la PR può essere mergiata.
\end{enumerate}

\subsection{Merge della Pull Request}

Esistono tre strategie di merge supportate da GitHub:

\begin{itemize}
  \item \textbf{Merge commit}: mantiene la storia di tutti i commit del branch;
  \item \textbf{Squash and merge}: condensa tutti i commit del branch in un singolo commit su \texttt{main};
  \item \textbf{Rebase and merge}: applica i commit del branch in cima a \texttt{main} senza commit di merge.
\end{itemize}

Per il progetto QuantSys Institutional si raccomanda \textbf{Squash and merge} per la maggior parte delle pull request, in modo da mantenere la storia di \texttt{main} pulita e leggibile.

\subsubsection{Step finale: cancellare il branch}

Dopo il merge, GitHub propone di cancellare il branch remoto. Si consiglia di farlo per mantenere la lista dei branch ordinata. Anche il branch locale può essere cancellato:

\begin{lstlisting}
git checkout main
git pull origin main
git branch -d feature/42-nome-descrittivo
\end{lstlisting}
